\chapter{المتنوعات الجزئية في \texorpdfstring{$\R^n$}{Rn}}
\label{ch:b3:submanifolds}

الكرات، والحلقات الدائرية، وزمر الدوران: ليست الموائل الطبيعية
للهندسة والميكانيكا فضاءاتٍ متجهية، بل \emph{مجموعاتٍ منحنية
تبدو مسطّحة عن قرب}. ويعطي هذا الفصل تلك العبارة معنًى دقيقًا
--- أي المتنوعات الجزئية في $\R^n$ --- ويعطي حسابها التفاضلي.
والأساس هو مبرهنة الدالة العكسية، مبرهَنًا عليها هنا بنقطة باناخ
الصامدة؛ وكل ما عداها تغييرُ إحداثيات: أي الأوصاف الأربعة
المتكافئة \omterm{thm:b3:submanifolds:characterizations}{للمتنوعة الجزئية} (التسوية المحلية، ومجموعات
المستويات، والمنحنيات البيانية، والوسطنات)، والفضاءات المماسّة،
والأمثلة المقيَّدة \omterm{thm:b3:submanifolds:lagrange}{بمضاعفات لاغرانج} --- التي تعيد، بوصفها
عرضًا وداعيًّا، البرهانَ على المبرهنة الطيفية للمصفوفات
المتناظرة في ثلاثة أسطر من الهندسة. وتبني مسألة نهاية الأسبوع
زمرة الدوران $SO(3)$ وغطاءها المزدوج الرباعي: أي الجبر (وقد
كبرت $Q_8$ في \cref{ch:b3:groups}) يلاقي الهندسة.

\section{مبرهنة الدالة العكسية}

\begin{theorem}[مبرهنة الدالة العكسية]
\label{thm:b3:submanifolds:ift}
لتكن $U \subseteq \R^n$ مفتوحة، ولتكن $f \colon U \to \R^n$ من الصنف $\mathcal C^1$، وليكن
$a \in U$ بحيث يكون $Df(a)$ قابلًا للقلب. عندئذٍ توجد مجموعتان
مفتوحتان $V \ni a$ و $W \ni f(a)$ بحيث يكون $f \colon V \to W$ تقابلًا
بمقلوب من الصنف $\mathcal C^1$، و\index{مبرهنة الدالة العكسية}
\[
D(f^{-1})(y) = \bigl(Df(f^{-1}(y))\bigr)^{-1}
\qquad (y \in W).
\]
وإذا كانت $f$ من الصنف $\mathcal C^k$، كان $f^{-1}$ كذلك.
\end{theorem}

\begin{proof}
نعيِّر: فباستبدال $f$ بالمقدار $x \mapsto Df(a)^{-1}\bigl(f(a +
x) - f(a)\bigr)$، يمكننا أن نفترض $a = 0$
و $f(0) = 0$ و $Df(0) =
I$ (وتنتج العبارة العامة بالتركيب مع
التقابلات التآلفية). ونكتب $f(x) = x + g(x)$: فإن $Dg(0) = 0$، وباتصال
$Dg$ نختار $r > 0$ بحيث $\vertiii{Dg(x)} \leq \frac12$ على $\bar B(0, r)$؛ وتعطي
متراجحة القيم المتوسطة أن $\norm{g(x) - g(x')} \leq \frac12
\norm{x - x'}$ هناك.

\emph{التقابل على \omterm{def:b3:topology:topology}{جوار}.} من أجل $y \in B(0,
\frac r2)$، يعني حل $f(x) = y$
إيجادَ نقطة صامدة للتطبيق $\Phi_y(x) = y - g(x)$؛ ويرسل $\Phi_y$ الكرة
$\bar B(0, r)$ في نفسها ($\norm{\Phi_y(x)} \leq \norm y + \frac12\norm x \leq
r$) وهو ليبشيتزي بالثابت
$\frac12$: فتعطي باناخ (\cref{thm:b3:complete:banach}) حلًّا وحيدًا $x = \varphi(y) \in \bar B(0,r)$. وفوق
ذلك تكون $f$ متباينة على $\bar B(0,r)$:
\[
\norm{f(x) - f(x')} \geq \norm{x - x'} - \norm{g(x) - g(x')}
\geq \tfrac12\norm{x - x'} .
\tag{$*$}
\]
ونضع $W = B(0, \frac r2)$ و $V = f^{-1}(W)\cap B(0, r)$: فهما مفتوحتان (بالاتصال)، مع $f \colon V \to W$
تقابليًّا.

\emph{اتصال المقلوب وقابليته للاشتقاق.} يقول $(*)$ إن $\varphi = f^{-1}$
ليبشيتزي بالثابت $2$. ونثبّت $y_0 = f(x_0)
\in W$؛ فتعطي قابلية $A = Df(x_0)$
للقلب (إذ مسافته إلى $I$ هي $\leq \frac12$: بنويمان،
\cref{prop:b3:banach:neumann}) وقابليةُ $f$ للاشتقاق، من أجل $y = f(x)$ بجوار
$y_0$:
\[
\varphi(y) - \varphi(y_0) - A^{-1}(y - y_0)
= -A^{-1}\bigl(f(x) - f(x_0) - A(x - x_0)\bigr)
= -A^{-1}\,o(\norm{x - x_0}) = o(\norm{y - y_0}),
\]
باستعمال $(*)$ لتحويل $\norm{x - x_0} \leq 2\norm{y - y_0}$: فتكون $\varphi$ قابلة للاشتقاق
عند $y_0$ بالتفاضل المقلوب. وأما اتصال $y \mapsto D\varphi(y) =
Df(\varphi(y))^{-1}$: فبتركيب تطبيقات
متصلة (إذ القلب متصل، \cref{prop:b3:banach:neumann}): أي $\varphi \in \mathcal C^1$؛ وترقيةُ الصيغة
نفسها تعطي $\mathcal C^k$.
\end{proof}

\begin{theorem}[مبرهنة الدوال الضمنية]
\label{thm:b3:submanifolds:implicit}
لتكن $F \colon U \subseteq \R^p\times\R^q \to \R^q$ من الصنف $\mathcal C^1$ بجوار $(a, b)$، مع
$F(a,b) = 0$، ولنفترض أن التفاضل الجزئي $D_yF(a,b) \in \mathcal L(\R^q)$ قابل للقلب.
عندئذٍ توجد جواران $A \ni a$ و $B \ni b$ وتطبيقٌ $\psi \colon A \to B$ من
الصنف $\mathcal C^1$ يحقق
\[
\bigl\{(x, y)\in A\times B : F(x,y) = 0\bigr\}
= \{(x, \psi(x)) : x \in A\},
\]
و $D\psi(x) = -D_yF(x, \psi(x))^{-1}\,D_xF(x,
\psi(x))$.\index{مبرهنة الدوال الضمنية}
\end{theorem}

\begin{proof}
نطبّق \cref{thm:b3:submanifolds:ift} على $\Theta(x, y) = (x,
F(x,y))$: فتفاضله عند $(a,b)$، وهو كتلي
مثلثي بكتلتين قطريتين قابلتين للقلب $I$ و $D_yF$، قابلٌ للقلب.
وللمقلوب المحلي الشكلُ $\Theta^{-1}(x, z) = (x, h(x,
z))$؛ ونضع $\psi(x) = h(x, 0)$: عندئذٍ يكون
$F(x, y) = 0$ إذا وفقط إذا كان $\Theta(x,y) = (x, 0)$، إذا وفقط إذا كان
$y = \psi(x)$، محليًّا. وأما الصيغة: فباشتقاق $F(x, \psi(x)) = 0$ بقاعدة
السلسلة.
\end{proof}

\section{المتنوعات الجزئية: أربعة تعاريف}

\begin{theorem}[التمييزات المتكافئة]
\label{thm:b3:submanifolds:characterizations}
لتكن $M \subseteq \R^n$، و $d \in \{0, \dots, n\}$، و $k \geq
1$. تتكافأ العبارات التالية، من
أجل كل نقطة $a \in M$ (وتكون $M$ \emph{متنوعةً جزئية ذات بُعد
$d$ من الصنف $\mathcal C^k$}\index{متنوعة جزئية} إذا صحّت عند
كل $a
\in M$):
\begin{enumerate}
  \item (\emph{التسوية}) يوجد تماثل تفاضلي $\Phi$ من الصنف
        $\mathcal C^k$ من مفتوحة $\Omega \ni a$ على مفتوحة
        $\Omega' \subseteq \R^n$ يحقق
        \[
        \Phi(M\cap\Omega) = \Omega' \cap
        \bigl(\R^d\times\{0\}\bigr).
        \]
  \item (\emph{مجموعة المستوى}) يوجد \emph{غمرٌ غامر}
        $F \colon \Omega \to \R^{n-d}$ من الصنف $\mathcal C^k$ (أي إن $DF(x)$ غامر)
        على مفتوحة $\Omega \ni a$ يحقق $M\cap\Omega = F^{-1}(0)$.
  \item (\emph{المنحني البياني}) إلى غاية تبديل الإحداثيات،
        تكون $M$ محليًّا المنحنيَ البياني لتطبيق $\psi
        \colon A \subseteq \R^d \to \R^{n-d}$ من
        الصنف $\mathcal C^k$.
  \item (\emph{الوسطنة}) يوجد \emph{غمرٌ} $\varphi \colon A \subseteq \R^d \to
        \R^n$ من الصنف
        $\mathcal C^k$ (أي إن $D\varphi(u)$ متباين) مع $A$
        مفتوحة، و $\varphi$ تماثلٌ طوبولوجي من $A$ على $M \cap
        \Omega$
        من أجل مفتوحة $\Omega \ni a$ ما.
\end{enumerate}
\end{theorem}

\begin{proof}
(1)$\Rightarrow$(2): $F = (\Phi_{d+1}, \dots, \Phi_n)$ (أي الإحداثيات الأخيرة للتطبيق
$\Phi$): وهو غمرٌ غامر (لأن $D\Phi$ قابل للقلب).
(2)$\Rightarrow$(3): $DF(a)$ غامر: ومنه يكون محدّدٌ جزئي ما من
الرتبة $q \times q$ للمصفوفة الياكوبية قابلًا للقلب (حيث
$q = n - d$)؛ وبعد تبديل الإحداثيات، يكون $D_yF(a)$ قابلًا
للقلب، وتعبّر مبرهنة الدوال الضمنية (\cref{thm:b3:submanifolds:implicit}) عن $M$
محليًّا بوصفها منحنيًا بيانيًّا $y = \psi(x)$.
(3)$\Rightarrow$(4): $\varphi(x) = (x, \psi(x))$: وهو غمرٌ (إذ تفاضله $\bigl(\begin{smallmatrix}I\\
D\psi\end{smallmatrix}\bigr)$ متباين)،
وتماثلٌ طوبولوجي على المنحني البياني (ومقلوبه: الإسقاط، وهو
متصل).
(4)$\Rightarrow$(1): لتكن $\varphi(u_0) = a$؛ نكمّل $\operatorname{im}D\varphi(u_0)$ بمتمم $E$ (حيث
$\dim E = n - d$) ونعرّف $\Theta(u, v) = \varphi(u) + v$ على $A\times E$: فيكون $D\Theta(u_0, 0)$
تقابليًّا (إذ تحتوي الصورة $\operatorname{im}D\varphi(u_0)$ و $E$)، ومنه يكون $\Theta$
تماثلًا تفاضليًّا محليًّا (\cref{thm:b3:submanifolds:ift})؛ ويسوّي مقلوبه $\Phi$:
أي إن نقاط $M$ بجوار $a$ هي بالضبط $\varphi(u) = \Theta(u, 0)$ --- ومن أجل ذلك،
تضمن فرضية \omterm{def:b3:topology:continuity}{التماثل الطوبولوجي} في (4) أن $M\cap\Omega$، من أجل
$\Omega$ صغيرة، لا تحتوي أي صفائح أخرى (إذ يجب استبعاد $\varphi(u') + v = m \in
M$
القريبة من $a$ مع $v \neq 0$ صغيرة: فإن $m = \varphi(u'')$ من أجل $u''$ ما
بجوار $u_0$ بخاصية \omterm{def:b3:topology:continuity}{التماثل الطوبولوجي}، ويفرض التباين المحلي
للتطبيق $\Theta$ أن $v = 0$). عندئذٍ $\Phi(M\cap\Omega) = (A\times\{0\})
\cap\Phi(\Omega)$ إلى غاية التقليص.
\end{proof}

\begin{example}\label{ex:b3:submanifolds:examples}
الكرة $S^{n-1} = \{\norm x^2 = 1\}$: مجموعة مستوى للغمر الغامر $F(x) = \norm x_2^2 - 1$ على $\R^n\setminus\{0\}$
($DF(x) = 2x^{\mathsf T} \ne 0$): أي \omterm{thm:b3:submanifolds:characterizations}{متنوعة جزئية} $\mathcal C^\infty$ بُعدها $n - 1$. والحلقة
الدائرية في $\R^3$: مجموعة مستوى للمقدار $\bigl(\sqrt{x^2 + y^2} -
R\bigr)^2 + z^2 - r^2$ (حيث
$0 < r < R$). والمخروط $\{x^2 + y^2 =
z^2\}$ \emph{ليس} \omterm{thm:b3:submanifolds:characterizations}{متنوعة جزئية} عند $0$
(\cref{exo:b3:submanifolds:1}). وزمر المصفوفات: $SL_n$ و $O_n$ متنوعتان جزئيتان
في $M_n(\R)$ (\cref{exo:b3:submanifolds:5,exo:b3:submanifolds:6}) --- وهي نقطة انطلاق نظرية لي.
\end{example}

\section{الفضاءات المماسّة}

\begin{definition}\label{def:b3:submanifolds:tangent}
لتكن $M$ متنوعةً جزئية ذات بُعد $d$ وليكن $a \in M$.
\emph{الفضاء المماسّ}\index{فضاء مماسّ} $T_aM$ هو مجموعة متجهات
السرعة $\gamma'(0)$ للمنحنيات $\gamma \colon
\intoo{-\varepsilon}\varepsilon \to M$ من الصنف $\mathcal C^1$
التي تحقق $\gamma(0) = a$.
\end{definition}

\begin{proposition}\label{prop:b3:submanifolds:tangent}
$T_aM$ فضاء متجهي جزئي من $\R^n$ بُعده $d$، و:
\begin{enumerate}
  \item إذا كان $M = F^{-1}(0)$ محليًّا مع $F$ غمرًا غامرًا:
        $T_aM = \ker DF(a)$؛
  \item وإذا كانت $M$ موسطَنة بالغمر $\varphi$ ($\varphi(u_0) = a$):
        $T_aM =
        \operatorname{im}D\varphi(u_0)$.
\end{enumerate}
\end{proposition}

\begin{proof}
تحقق المنحنيات في $M$ أن $F(\gamma(t)) = 0$؛ وتعطي قاعدة السلسلة عند $0$ أن
$DF(a)\gamma'(0) = 0$: أي $T_aM \subseteq \ker DF(a)$. وبالعكس، تنقل التسوية (\cref{thm:b3:submanifolds:%
characterizations}(1))
مستقيماتِ $\R^d\times\{0\}$ إلى منحنيات من $M$: فتتحقق كل متجهة من فضاء
جزئي ذي بُعد $d$؛ وبمقارنة الأبعاد ($\dim\ker DF(a) = n - (n - d) = d$) يُفرض التساوي في
(1)، وتعطي حجة النقل نفسها (2) (بتطبيق $D\varphi(u_0)$ على
المستقيمات في $A$؛ وبالأبعاد مرة أخرى).
\end{proof}

\begin{theorem}[مضاعفات لاغرانج]
\label{thm:b3:submanifolds:lagrange}
ليكن $M = F^{-1}(0)$ مع $F = (F_1, \dots, F_q) \colon \Omega
\to \R^q$ غمرًا غامرًا من الصنف
$\mathcal C^1$، ولتكن $f \colon
\Omega \to \R$ من الصنف $\mathcal C^1$. إذا كان
للتضييق $f\restriction_M$ قيمةٌ قصوى محلية عند $a \in M$، فتوجد أعداد
حقيقية وحيدة $\lambda_1, \dots, \lambda_q$ (أي \emph{\omterm{thm:b3:submanifolds:lagrange}{مضاعفات لاغرانج}})
تحقق\index{مضاعفات لاغرانج}
\[
\nabla f(a) = \lambda_1\nabla F_1(a) + \dots +
\lambda_q\nabla F_q(a) .
\]
\end{theorem}

\begin{proof}
من أجل كل منحنٍ $\gamma$ في $M$ يمرّ بالنقطة $a$: يكون للمقدار
$t \mapsto
f(\gamma(t))$ قيمةٌ قصوى محلية عند $0$، ومنه $0 =
\frac{\dd}{\dd t}f(\gamma(t))\big|_0 = \langle\nabla f(a),
\gamma'(0)\rangle$: أي $\nabla f(a) \perp T_aM = \ker DF(a)$
(\cref{prop:b3:submanifolds:tangent}). والآن $\ker DF(a)^\perp
= \operatorname{im}DF(a)^{\mathsf T}$ (أي متطابقة الرتبة أو التعامد في
السنة الجامعية 2، أو \cref{exo:b3:hilbert:8} في البُعد المنتهي: $(\ker T)^\perp = \operatorname{im}T^*$)،
وهو مولَّد بالتدرّجات $\nabla F_i(a)$ --- وهي مستقلة لأن
$DF(a)$ غامر: فتوجد المضاعفات وتكون وحيدة.
\end{proof}

\begin{example}[المبرهنة الطيفية، هندسيًّا]
\label{ex:b3:submanifolds:rayleigh}
لتكن $A$ مصفوفةً حقيقية متناظرة من الرتبة $n\times n$ ولنعظّم
$f(x) = \langle Ax, x\rangle$ على الكرة $S^{n-1}$ (وهي متراصة: فتُبلَغ القيمة العظمى
عند $v_1$ ما). وبلاغرانج مع $F(x) = \norm x^2 - 1$: يعطي $\nabla f = 2Ax$
و $\nabla F =
2x$ أن $Av_1 = \lambda_1v_1$ --- أي متجهة ذاتية، مع $\lambda_1 = \max_{S^{n-1}}\langle Ax, x\rangle$. ثم نضيّق
$A$ على $v_1^\perp$ (وهو صامد: $\langle Av, v_1\rangle =
\langle v, Av_1\rangle = \lambda_1\langle v, v_1\rangle = 0$) ونكرّر: فنجد أساسًا
متعامدًا متجانسًا من المتجهات الذاتية. وهي المبرهنة الطيفية في
السنة الجامعية 2، معادًا البرهان عليها بالأمثلة المحضة --- وظلّ
الحجة نفسها في البُعد اللانهائي برهن على \cref{lem:b3:spectral:existence}.
\end{example}

\begin{method}\label{met:b3:submanifolds:method}
لتبرهن على أن مجموعةً متنوعةٌ جزئية: أبرزها محليًّا بوصفها
$F^{-1}(0)$ مع $DF$ غامرًا \emph{على المجموعة} (وهو الطريق
الأشيع)، أو بوصفها منحنيًا بيانيًّا. ولتحسب بُعدها وفضاءها
المماسّ: $d = n - q$ و $T_a = \ker DF(a)$. ولتُمثِّل عليها: لاغرانج ---
وتحقق دائمًا من التراص (أو القسرية) أولًا، بحيث توجد قيمة قصوى
يمكن للمبرهنة أن تنطبق عليها، وتذكّر أن معادلة المضاعفات
\emph{لازمة} فحسب: فاجمع كل النقاط الحرجة، ثم قارن القيم. وأما
زمر المصفوفات، فاشتقّ المنحنيات عند المطابقة لتعيّن الفضاءات
المماسّة.
\end{method}

\section{تمارين}

\begin{exercise}[$\star$]\label{exo:b3:submanifolds:1}
(a) تحقق من أن ما يلي متنوعات جزئية $\mathcal C^\infty$ وأعطِ أبعادها:
$S^{n-1}$؛ والقطع الزائد $\{x^2 + y^2 - z^2 = 1\}$؛ والحلقة الدائرية في
\cref{ex:b3:submanifolds:examples}.
(b) برهن على أن المخروط $C = \{x^2 + y^2 = z^2\} \subseteq
\R^3$ ليس \omterm{thm:b3:submanifolds:characterizations}{متنوعة جزئية} ذات بُعد $2$
عند $0$: عيّن عدد المركّبات المترابطة للمجموعة $\bigl(C\setminus\{0\}\bigr)\cap B(0,\varepsilon)$، وقارنه
بالعدد الذي كانت التسوية (\cref{thm:b3:submanifolds:characterizations}(1)) ستفرضه على مستوٍ
منزوعةً منه نقطة.
\end{exercise}

\begin{exercise}[$\star$]\label{exo:b3:submanifolds:2}
احسب الفضاءات المماسّة:
(a) $T_aS^{n-1}$ من أجل أي $a$ (والجواب: $a^\perp$)؛
(b) المستوي المماسّ للحلقة الدائرية في \cref{ex:b3:submanifolds:examples} عند نقطة
كيفية من خط الاستواء الخارجي $\{z = 0,\ x^2 + y^2 = (R + r)^2\}$؛
(c) المستقيم المماسّ للولب $\varphi(t) = (\cos t,
\sin t, t)$ عند $\varphi(t_0)$، متحققًا
من \cref{prop:b3:submanifolds:tangent}(2).
\end{exercise}

\begin{exercise}[$\star\star$]\label{exo:b3:submanifolds:3}
لتكن $f(x, y) = (x^2 - y^2,\ 2xy)$ (أي $z \mapsto z^2$).
(a) عند أي النقاط تنطبق \cref{thm:b3:submanifolds:ift}؟
(b) برهن على أن $f$ قابلة للقلب محليًّا لا شاملًا على $\R^2\setminus\{0\}$،
وأبرز صراحةً المقلوبين المحليين المعرَّفين على \omterm{def:b3:topology:topology}{جوار} للنقطة
$(1, 0)$ (أي فرعَي الجذر التربيعي).
(c) والمناقشة نفسها من أجل تطبيق \omterm{ex:b3:product:polar}{الإحداثيات القطبية} $(r,
\theta) \mapsto (r\cos\theta, r\sin\theta)$.
\end{exercise}

\begin{exercise}[$\star\star$]\label{exo:b3:submanifolds:4}
(ورقة ديكارت) لتكن $F(x,y) = x^3 + y^3 - 3xy$ و $\mathcal C =
F^{-1}(0)$.
(a) برهن على أنه بجوار كل نقطة من $\mathcal C$ غير المبدأ،
تكون $\mathcal C$ متنوعةً جزئية ذات بُعد $1$، وهي محليًّا منحنٍ
بياني في $x$ أو في $y$ (فأيهما، وأين؟).
(b) احسب المستقيم المماسّ عند $(\frac32, \frac32)$.
(c) وماذا يحدث عند المبدأ؟ (إذ يتقاطع فرعان: أبرز منحنيين من
الصنف $\mathcal C^1$ في $\mathcal C$ يمرّان بالنقطة $0$
بسرعتين مستقلتين، واخلص إلى أنه لا توجد أي تسوية.)
\end{exercise}

\begin{exercise}[$\star\star$]\label{exo:b3:submanifolds:5}
لتكن $F(M) = M^{\mathsf T}M$ من $M_n(\R)$ إلى فضاء المصفوفات المتناظرة $S_n$.
(a) برهن على $DF(M)(H) = M^{\mathsf T}H + H^{\mathsf T}M$ وعلى أن $DF(M)$ غامر على $S_n$ عند كل
$M \in O_n$ \emph{(بإعطاء $S \in S_n$، جرّب $H = \frac12MS$)}.
(b) اخلص إلى أن $O_n = F^{-1}(I)$ \omterm{thm:b3:submanifolds:characterizations}{متنوعة جزئية} متراصة $\mathcal C^\infty$ بُعدها
$\frac{n(n-1)}2$، مع $T_IO_n = \{H : H^{\mathsf T} = -H\}$، أي المصفوفات المتخالفة.
(c) برهن على $\eu^{tH} \in O_n$ من أجل كل $H$ متخالفة: أي إن الاتجاهات
المماسّة تتكامل إلى منحنيات في الزمرة.
\end{exercise}

\begin{exercise}[$\star\star$]\label{exo:b3:submanifolds:6}
(a) برهن على أن للمقدار $\det \colon M_n(\R) \to \R$ التفاضلَ $D\det(M)(H) = \operatorname{tr}\bigl(\operatorname{com}(M)
^{\mathsf T}H\bigr)$، وهو غير معدوم
عند كل $M \in SL_n$.
(b) اخلص إلى أن $SL_n(\R)$ \omterm{thm:b3:submanifolds:characterizations}{متنوعة جزئية} بُعدها $n^2 - 1$ مع
$T_ISL_n = \{H : \operatorname{tr}H = 0\}$.
(c) وهل $GL_n(\R)$ \omterm{thm:b3:submanifolds:characterizations}{متنوعة جزئية}؟ وما بُعدها؟
\end{exercise}

\begin{exercise}[$\star\star$]\label{exo:b3:submanifolds:7}
\omterm{thm:b3:submanifolds:lagrange}{بمضاعفات لاغرانج}:
(a) جد القيم القصوى للدالة $f(x,y) = xy$ على الدائرة $x^2 +
y^2 = 1$؛
(b) برهن على أنه من بين جميع المتجهات الاحتمالية $(p_1, \dots,
p_n)$
(الموجبة التي مجموعها $1$)، تُعظَّم الإنتروبيا $-\sum
p_i\ln p_i$ عند
التوزيع المنتظم بالضبط؛
(c) جد نقطة القطع الناقص $\{x^2/4 + y^2 = 1\}$ الأقرب إلى $(1, 0)$، وتحقق من
معادلة المضاعفات هندسيًّا (باصطفاف الناظم).
\end{exercise}

\begin{exercise}[$\star\star\star$]\label{exo:b3:submanifolds:8}
اكتب \cref{ex:b3:submanifolds:rayleigh} كاملًا: برهن بالتراجع على أن مصفوفةً حقيقية
متناظرة تقبل أساسًا متعامدًا متجانسًا من المتجهات الذاتية، مع
$\lambda_1 \geq \dots \geq \lambda_n$ القيمَ القصوى المقيَّدة المتتالية لنسبة رايلي. ثم
استنتج صيغ كورانت--فيشر في \cref{exo:b3:spectral:8} في البُعد المنتهي
مباشرةً من هذا البناء.
\end{exercise}

\begin{exercise}[$\star\star\star$]\label{exo:b3:submanifolds:9}
(متراجحة هادامار) من أجل $M \in GL_n(\R)$ ذي الأعمدة $c_1, \dots, c_n$:
\[
\abs{\det M} \;\leq\; \prod_{i=1}^n\norm{c_i}_2 ,
\]
مع التساوي إذا وفقط إذا كانت الأعمدة متعامدة. \emph{(أرجع إلى
أعمدة معيارها $1$ بالتحجيم؛ وعظّم $\det$ على الجداء المتراص
للكرات $(S^{n-1})^n$؛ وعند معظِّم، تعطي لاغرانج في كل عمود على
حدة $\nabla_{c_i}\det =
\lambda_ic_i$، ويكون $\nabla_{c_i}\det$ العمودَ ذا الرتبة $i$ من $\operatorname{com}(M)$:
فاستنتج أن $M^{\mathsf T}M$ قطرية، ومنه $= I$، ومنه
$\det = \pm1$.)} والقراءة الهندسية: أن حجم متوازي السطوح لا
يتجاوز جداء أطوال حروفه.
\end{exercise}

\begin{exercise}[$\star\star$]\label{exo:b3:submanifolds:10}
بجوار أي من نقاطها تكون الدائرة $S^1$ منحنيًا بيانيًّا $y =
\psi(x)$؟
ومنحنيًا بيانيًّا $x = \chi(y)$؟ تحقق من تمييز المنحني البياني
(\cref{thm:b3:submanifolds:%
characterizations}(3)) صراحةً عند $(1, 0)$، واشرح في جملة واحدة لماذا
يكفي دائمًا تبديلٌ \emph{ما} للإحداثيات لكن لا يصلح تبديل واحد
دائمًا.
\end{exercise}

\begin{exercise}[$\star\star$]\label{exo:b3:submanifolds:11}
(الزمرة المتعامدة بوصفها \omterm{thm:b3:submanifolds:characterizations}{متنوعة جزئية}، كمّيًّا)
(a) برهن على أن $O_n = \{M : M^{\mathsf T}M = I\}$ متراصة: محدودة (إذ كل عمود متجهة
واحدية، ومنه $\norm M \leq
\sqrt n$ بالمعيار المصفوفي الإقليدي) ومغلقة.
(b) برهن على أن فضاءها المماسّ عند $I$ هو فضاء المصفوفات
المتخالفة، وبُعده $\frac{n(n-1)}2$، وعند $A \in O_n$ عام:
$T_AO_n = \{AK : K^{\mathsf T} =
-K\}$.
(c) استنتج أن التطبيق $t \mapsto A\exp(tK)$ هو، من أجل كل $K$ متخالفة، منحنٍ
في $O_n$ يمرّ بالنقطة $A$ بسرعة $AK$ \emph{(تحقق من $\exp(tK) \in O_n$
باستعمال $\exp(X)^{\mathsf T} = \exp(X^{\mathsf T})$ و $\exp(-X)\exp(X) = I$)}: فتتحقق كل متجهة مماسّة بمنحنٍ
صريح، دون حاجة إلى مبرهنة الدوال الضمنية.
\end{exercise}

\begin{exercise}[$\star\star$]\label{exo:b3:submanifolds:12}
(النقاط الحرجة للمسافة) لتكن $M \subseteq \R^n$ متنوعةً جزئية وليكن
$p \notin M$. برهن على أنه إذا صغّرت $x_0 \in M$ المسافة إلى
$p$ (وتوجد نقطة كهذه حين تكون $M$ مغلقة وغير خالية --- ولماذا؟)،
فإن
\[
p - x_0 \;\perp\; T_{x_0}M
\]
\emph{(باشتقاق $t \mapsto \norm{\gamma(t) - p}^2$ على امتداد منحنيات في $M$)}. واستنتج: أن
أقرب نقطة على كرة تقع على نصف المستقيم المارّ بالمركز؛ واستعمل
الشرط لحساب المسافة من $p = (2, 0)$ إلى القطع المكافئ
$y = x^2$ (بالإرجاع إلى معادلة من الدرجة الثالثة وحلّها عدديًّا
بثلاثة أرقام).
\end{exercise}

\section{مسألة: \texorpdfstring{$SO(3)$}{SO(3)} والرباعيات}

\begin{problem}[{مسألة نهاية الأسبوع --- الدورانات، وزمرة
$S^3$، والغطاء المزدوج}]\label{pb:b3:submanifolds:1}
الرباعيات $\mathbb H = \{t + x\mathrm i + y\mathrm j +
z\mathrm k\}$ --- أي الجبر الذي تحتوي زمرةُ عناصره القابلة
للقلب $Q_8$ في \cref{pb:b3:groups:1} --- توسطن الدورانات الثلاثية الأبعاد
مرتين: فتطبيق «الاقتران برباعية واحدية» تشاكلٌ غامر $S^3 \to
SO(3)$
نواته $\{\pm1\}$. ونبني كل شيء. تذكيرًا أو تعريفًا: الضرب ثنائي
الخطية بالمعنى $\R$ مع $\mathrm i^2 = \mathrm j^2 = \mathrm k^2 =
\mathrm{ijk} = -1$؛ ومرافق $q = t + x\mathrm i +
y\mathrm j + z\mathrm k$ هو $\bar q = t - x\mathrm i -
y\mathrm j - z\mathrm k$؛
و $N(q) = q\bar q = t^2 + x^2 + y^2
+ z^2$.

\medskip
\noindent\textbf{الجزء الأول --- الجبر $\mathbb H$ والزمرة
$S^3$.}
\begin{enumerate}
  \item تحقق من أن $\mathbb H$ جبرٌ تجميعي على $\R$ مركزه
        $\R$، ومن أن $\overline{pq} =
        \bar q\,\bar p$، ومن أن $N(pq) = N(p)N(q)$ \emph{(وطريق نظيف
        واحد: مثّل $q$ بالمصفوفة العقدية من الرتبة $2\times2$
        $\bigl(\begin{smallmatrix}\alpha & \beta\\
        -\bar\beta & \bar\alpha\end{smallmatrix}\bigr)$، $q =
        \alpha + \beta\mathrm j$، واستعمل $\det$)}.
  \item استنتج أن كل $q \neq 0$ قابل للقلب ($q^{-1} =
        \bar q/N(q)$): فيكون
        $\mathbb H$ حقلًا (غير تبديلي)، وتكون $S^3 = \{N(q) = 1\}$ زمرةً ---
        ومتنوعةً جزئية متراصة ذات بُعد $3$ في $\R^4$
        (\cref{ex:b3:submanifolds:examples}).
\end{enumerate}

\medskip
\noindent\textbf{الجزء الثاني --- تشاكل الدوران.} نطابق $\R^3$
مع \emph{الرباعيات الصرفة} $P = \{x\mathrm i +
y\mathrm j + z\mathrm k\}$، ونعرّف من أجل $q \in S^3$
المقدارَ $\rho_q(v) = q\,v\,\bar q$.
\begin{enumerate}[resume]
  \item برهن على أن $\rho_q$ يرسل $P$ في $P$ \emph{(إذ
        الرباعيات الصرفة هي التي تحقق $\bar v = -v$)}، وأنه
        خطي بالمعنى $\R$، ويحفظ المعيار، وأن $\rho
        \colon q \mapsto \rho_q$ تشاكل زمر
        $S^3
        \to O(3)$.
  \item احسب النواة: $\rho_q = \mathrm{id}$ إذا وفقط إذا كان $q$ يتبادل مع
        $\mathrm i, \mathrm j, \mathrm k$، إذا وفقط إذا كان $q \in \R\cap S^3 = \{\pm1\}$.
  \item اكتب $q = \cos\frac\theta2 + \sin\frac\theta2\,u$ مع $u \in P$ و $N(u) = 1$ (ولماذا يكون
        هذا ممكنًا دائمًا من أجل $q \in S^3$؟). برهن على أن
        $\rho_q$ يثبّت $u$، وأنه على المستوي $u^\perp\cap P$
        يفعل بوصفه الدورانَ بالزاوية $\theta$ \emph{(احسب
        $\rho_q(w)$ من أجل $w \perp u$ باستعمال $uw = -wu$
        من أجل الوحدات الصرفة المتعامدة --- وبرهن على هذه
        المتطابقة من جدول الضرب، أو من $uw + wu =
        -2\langle u, w\rangle$)}.
  \item اخلص: أن $\operatorname{im}\rho \subseteq SO(3)$ (إذ كل $\rho_q$ دورانٌ بمحور وزاوية
        كما حُسبا --- ومحدده $+1$ باتصال $q
        \mapsto \det\rho_q$ على $S^3$
        المترابطة، أو مباشرةً)، وأن $\rho$ \emph{غامر} على
        $SO(3)$: أي إن لكل دوران في $\R^3$ محورًا (برهن: أن
        لمصفوفة حقيقية متعامدة من الرتبة $3\times3$ تحقق
        $\det = 1$ القيمةَ الذاتية $1$ --- بالنظر في كثير الحدود
        المميِّز) ومن ثَمّ يكون $\rho_q$ ما. والخلاصة:
        \[
        SO(3) \;\cong\; S^3/\{\pm 1\} .
        \]
\end{enumerate}

\medskip
\noindent\textbf{الجزء الثالث --- $SO(3)$ بوصفها \omterm{thm:b3:submanifolds:characterizations}{متنوعة
جزئية}؛ رودريغ.}
\begin{enumerate}[resume]
  \item برهن على أن $SO(3)$ \omterm{thm:b3:submanifolds:characterizations}{متنوعة جزئية} متراصة ذات بُعد $3$
        في $M_3(\R)$ مع $T_ISO(3) =$
        المصفوفات المتخالفة
        (\cref{exo:b3:submanifolds:5}؛ ويختار شرط المحدد اتحادَ مركّبات).
  \item من أجل المصفوفة المتخالفة $A_u$ المرافقة للمتجهة
        $u \in \R^3$ ($A_uv = u\wedge v$، أي الجداء المتجهي)، برهن على
        \emph{صيغة رودريغ}:
        \[
        \eu^{\theta A_u} = I + \sin\theta\,A_u + (1 -
        \cos\theta)\,A_u^2 \qquad (\norm u = 1)
        \]
        \emph{(انطلاقًا من $A_u^3 = -A_u$: اشطر المتسلسلة
        الأُسّية على امتداد قوى $A_u$)}، وعيّنها بوصفها
        الدورانَ بالمحور $u$ والزاوية $\theta$. واستنتج أن
        $\exp$ يرسل المصفوفات المتخالفة \emph{على} $SO(3)$.
  \item اربط بين الوسطنتين: برهن على أن $t
        \mapsto \rho_{q(t)}$ مع $q(t) = \cos\frac t2 +
        \sin\frac t2\,u$
        زمرةٌ ذات وسيط واحد من الدورانات مشتقتها عند $t = 0$
        هي $A_u$ --- فالأُسّيان الرباعي والمصفوفي يرويان
        القصة نفسها بنصف السرعة وبالسرعة الكاملة على التوالي.
\end{enumerate}

\medskip
\noindent\textbf{الجزء الرابع --- الغطاء المزدوج، محسوسًا.}
\begin{enumerate}[resume]
  \item برهن على أن المسار $q(t) = \cos\frac t2 + \sin\frac
        t2\,\mathrm k$، حيث $t \in \intcc0{2\pi}$، حلقةٌ في
        $SO(3)$ (إذ ترجع صورته $\rho_{q(t)}$ إلى المطابقة)
        لكن رفعه الرباعي \emph{ليس} حلقة: $q(2\pi) = -q(0)$. وبالمتابعة
        إلى $t = 4\pi$ ينغلق الرفع. واشرح في فقرة قصيرة ما
        يقوله ذلك: أن دورانًا بمقدار $2\pi$ غير قابل للفك
        اتصاليًّا بينما دورانًا بمقدار $4\pi$ قابل لذلك (خدعة
        الحزام)، لأن حلقات $SO(3)$ تُكشف في غطائها المزدوج
        $S^3$.
  \item واستنتج أيضًا العائد العملي: أن تركيب الدورانات
        $=$ ضرب الرباعيات (أي ما يعادل $4$ عمليات ضرب من
        المعطيات بدل $9$، ودون انجراف عن التعامد) --- وتحقق
        من ذلك على تركيب ربعَي دورة حول $\mathrm i$
        و $\mathrm j$: احسب محور الجداء وزاويته.
\end{enumerate}

\medskip
\noindent\textbf{الجزء الخامس --- المصفوفة الصريحة:
أولير--رودريغ.} نكتب $q = a + b\mathrm i + c\mathrm j +
d\mathrm k \in S^3$، بحيث $a^2 + b^2 + c^2 + d^2 = 1$.
\begin{enumerate}[resume]
  \item احسب $\rho_q(\mathrm i)$ كاملًا من جدول الضرب؛ ثم احصل على $\rho_q(\mathrm
        j)$
        و $\rho_q(\mathrm k)$ بالتعويض الدوري $\mathrm i \to \mathrm j \to \mathrm k
        \to \mathrm i$، $(b, c, d) \to (c, d, b)$ \emph{(وبرّر
        ذلك: أن تدوير $\mathrm i, \mathrm j,
        \mathrm k$ يمتدّ إلى تشاكل ذاتي للجبر
        $\mathbb
        H$، لأن العلاقات المعرِّفة متناظرة دوريًّا)}.
        واخلص إلى أن مصفوفة $\rho_q$ في الأساس $(\mathrm i, \mathrm j, \mathrm k)$ هي
        \emph{مصفوفة أولير--رودريغ}
        \[
        R_q = \begin{pmatrix}
        a^2 + b^2 - c^2 - d^2 & 2(bc - ad) & 2(bd + ac)\\
        2(bc + ad) & a^2 - b^2 + c^2 - d^2 & 2(cd - ab)\\
        2(bd - ac) & 2(cd + ab) & a^2 - b^2 - c^2 + d^2
        \end{pmatrix}.
        \]
  \item (قراءة دوران عكسيًّا) برهن على
        \[
        \operatorname{tr}R_q = 4a^2 - 1 = 1 + 2\cos\theta,
        \qquad
        \tfrac12\bigl(R_q - R_q^{\mathsf T}\bigr) =
        \sin\theta\,A_u,
        \]
        باصطلاح السؤالين 5 و8. واستنتج خوارزمية تستعيد $\pm q$
        من مصفوفة دوران $R$: الزاوية من الأثر؛ والمحور من
        الجزء المتخالف حين $0 < \theta < \pi$؛ وحين $\theta = \pi$،
        برهن على المتطابقة $R
        + I = 2\,uu^{\mathsf T}$ واستعملها.
  \item قيّم $R_q$ من أجل الجداء $q =
        \frac12(1 + \mathrm i + \mathrm j + \mathrm k)$ في السؤال 11: فتظهر
        مصفوفة تبديل. عيّن الدوران ووفّق بينه وبين المحور
        والزاوية المحسوبين في السؤال 11.
\end{enumerate}

\medskip
\noindent\textbf{الجزء السادس --- داخل $S^3$: $SU(2)$،
\omterm{ex:b3:groups:actions}{وأصناف الاقتران}، والأُسّيات.}
\begin{enumerate}[resume]
  \item برهن على أن \omterm{def:b3:representations:rep}{التمثيل} المصفوفي في السؤال 1 (ولنسمّه
        $\Phi$) يتضيّق إلى تماثل زمر من $S^3$ على \emph{الزمرة
        الواحدية الخاصة}
        \[
        SU(2) = \bigl\{U \in M_2(\C) : U^*U = I,\ \det U =
        1\bigr\}
        \]
        \emph{(من أجل الغمر، اكتب المعادلتين $U^{-1} = U^*$
        و $\det U = 1$ من أجل مصفوفة عقدية عامة من الرتبة
        $2\times2$)}.
  \item برهن على أن الجزء الحقيقي صامدٌ بالاقتران على $S^3$ ---
        أي $\operatorname{Re}(pq\bar p) =
        \operatorname{Re}q$ من أجل كل $p \in S^3$ --- وبالعكس، على أن
        رباعيتين واحديتين لهما الجزء الحقيقي نفسه مقترنتان في
        $S^3$ \emph{(بالإرجاع إلى نقل محور صرف واحدي على آخر،
        وهو ما يوفّره الجزء الثاني)}. وصف أصناف اقتران $S^3$
        هندسيًّا؛ وترجم ذلك إلى $SU(2)$ (أي مجموعات مستويات
        الأثر)؛ وأسقط بواسطة $\rho$: فيكون دورانان مقترنين في
        $SO(3)$ إذا وفقط إذا كانت لهما الزاوية نفسها $\theta \in \intcc0\pi$.
  \item عرّف $\exp$ على $\mathbb H$ بالمتسلسلة الأُسّية؛ وتحقق
        من التقارب المطلق، باستعمال $\abs{pq} = \abs p\,\abs q$ من أجل $\abs q =
        \sqrt{N(q)}$.
        وبرهن، من أجل $u$ صرفة واحدية ومن أجل $\theta
        \in \R$، على
        \[
        \exp(\theta u) = \cos\theta + \sin\theta\,u ,
        \]
        واستنتج أن $\exp$ يرسل الفضاء الفائق $P$ على $S^3$،
        وتحقق من أن $\rho_{\exp(su)} =
        \eu^{2sA_u}$: أي ظاهرة نصف الزاوية في السؤال 9
        من جديد.
  \item من أجل رباعيتين صرفتين $v, w$ برهن على قاعدة الجداء
        $vw = -\langle v, w\rangle + v\wedge w$، ومن ثَمّ على متطابقة \omterm{def:b3:groups:derived}{المبدِّل} $vw - wv = 2\,v\wedge w$؛ وبرهن
        أيضًا على $[A_v, A_w] = A_{v\wedge w}$ من أجل مصفوفات السؤال 8. واخلص إلى
        أن مشتقة $\rho$ عند $1$ على امتداد المنحنيات $t \mapsto \exp(tv)$
        هي التماثل الخطي $v \mapsto 2A_v$ من $P$ على
        المصفوفات المتخالفة، وأنه ينقل مبدِّل الرباعيات إلى
        مبدِّل المصفوفات.
\end{enumerate}

\medskip
\noindent\textbf{الجزء السابع --- البنية الشاملة.}
\begin{enumerate}[resume]
  \item (لا مقطع متصل) لنفترض أن $s \colon SO(3) \to
        S^3$ متصل مع $\rho \circ s =
        \operatorname{id}$. ومن
        أجل الحلقة $R(t) =
        \rho_{q(t)}$ في السؤال 10، ضع $\varepsilon(t) =
        s(R(t))\,q(t)^{-1}$ من أجل
        $t \in \intcc0{2\pi}$. برهن على أن $\varepsilon$ متصل بقيم في
        $\{\pm1\}$، واستنتج تناقضًا: أي إنه لا يوجد أي اختيار
        شامل متصل لرباعية واحدية تمثّل كل دوران.
  \item (نموذج الكرة) لتكن $\bar B \subseteq \R^3$ الكرةَ المغلقة ذات نصف
        القطر $\pi$ ولتكن $E(v) = \eu^{A_v}$، مع $E(0) = I$. برهن على أن
        $E$ يرسل $\bar B$ على $SO(3)$، وأنه متباين على الكرة
        المفتوحة، وأنه على كرة الحافة يطابق المتقابلين
        بالضبط: $E(\pi u) = E(-\pi u) = 2uu^{\mathsf T} - I$، دون أي تصادفات أخرى. ومن ثَمّ تكون
        $SO(3)$ هي الكرةَ وقد أُلصقت نقاط حافتها المتقابلة ---
        أي الفضاء الإسقاطي $\mathbb{RP}^3$ --- ويصير القطر
        حلقةَ السؤال 10 غير القابلة للتقلّص.
  \item برهن على أن $\rho_p\rho_q\rho_p^{-1} = \rho_{pq\bar
        p}$؛ وعلى أن قلوب $SO(3)$ (أي المصفوفات
        $R \neq
        I$ التي تحقق $R^2 = I$) هي بالضبط أنصاف الدورات
        $\rho_w$ حيث $w$ رباعية صرفة واحدية؛ وعلى أن مركز
        $SO(3)$ بديهي.
  \item برهن على أن كل دوران جداءٌ لنصفَي دورة: فمن أجل $q = \cos\frac\theta2 +
        \sin\frac\theta2\,u$،
        اختر رباعية صرفة واحدية $w \perp
        u$، وتحقق من أن $w' = qw$
        رباعية صرفة واحدية أيضًا، وتحقق من $\rho_q =
        \rho_{w'}\rho_w$. وأين يقع
        المحوران، وما الزاوية بينهما؟
  \item اخلص إلى الملخص الطوبولوجي: أن $SO(3)$ متراصة ومترابطة
        بالمسارات (بأعطِ برهانين: صورة متصلة للكرة $S^3$
        بالتطبيق $\rho$؛ وصورة $\exp$)، بينما تملك $O(3)$
        مركّبتين مترابطتين بالضبط، كلٌّ منهما مطابقة طوبولوجيًّا
        للزمرة $SO(3)$.

  \item (تركيب، بثلاث طرق) ليكن $R_1$ الدورانَ بالزاوية
        $\frac\pi2$ حول محور $z$ وليكن $R_2$ الدورانَ بالزاوية
        $\frac\pi2$ حول محور $x$. احسب محور $R_2R_1$ وزاويته:
        (أ) بضرب المصفوفتين من الرتبة $3\times3$ واستعمال
        الأثر والجزء المتخالف (الجزء الخامس)؛ (ب) بضرب
        الرباعيتين الواحديتين الموافقتين $q_2q_1$. وتحقق من
        توافق الجوابين: الزاوية $\frac{2\pi}3$، والمحور
        $\frac1{\sqrt3}(1, -1, 1)$.
  \item (\omterm{ex:b3:conformal:mobius}{تحويل كايلي}) من أجل $K$ متخالفة، برهن على أن $I + K$
        قابل للقلب وعلى أن
        \[
        C(K) = (I - K)(I + K)^{-1} \in SO(n),
        \]
        مع عدم كون $-1$ قيمةً ذاتية للمقدار $C(K)$ أبدًا؛
        وبرهن على أن $K \mapsto C(K)$ تقابل من المصفوفات
        المتخالفة على $\{R \in SO(n) : -1 \notin
        \operatorname{Sp}R\}$، بمقلوب $R \mapsto (I -
        R)(I + R)^{-1}$. (وهي خريطة
        ناطقة للزمرة $SO(n)$، رفيقةٌ للأُسّي $\exp$ المتسامي في
        \cref{exo:b3:submanifolds:11}.)
\end{enumerate}
\end{problem}
